\chapter{Configuración de Build}
\label{app:build}

\section*{Opciones de Compilación}

\begin{center}
\begin{tabular}{ll}
\toprule
\textbf{Flag} & \textbf{Efecto} \\
\midrule
\texttt{-DLISPY\_PYTHON=ON} & Habilita soporte Python (defecto) \\
\texttt{-DLISPY\_PYTHON=OFF} & Compila sin Python \\
\texttt{-DCMAKE\_BUILD\_TYPE=Release} & Optimización máxima \\
\texttt{-DCMAKE\_BUILD\_TYPE=Debug} & Símbolos de depuración \\
\bottomrule
\end{tabular}
\end{center}

\section*{Requisitos para Build con Python}

\begin{itemize}
\item CPython 3.14 development headers
\item CMake 3.20+
\item Ninja o Make
\item Compilador C99 (GCC, Clang, MSVC)
\end{itemize}

\section*{Build en Linux}

\begin{verbatim}
sudo apt install python3-dev cmake ninja-build
cmake -S . -B build -G Ninja -DLISPY_PYTHON=ON
ninja -C build
\end{verbatim}

\section*{Build en Windows}

\begin{verbatim}
cmake -S . -B build -G Ninja -DLISPY_PYTHON=ON
ninja -C build
\end{verbatim}

\section*{Build sin Python}

\begin{verbatim}
cmake -S . -B build -G Ninja -DLISPY_PYTHON=OFF
ninja -C build
\end{verbatim}

\section*{Estructura de Archivos Post-Build}

\begin{verbatim}
build/
├── bin/
│   └── lispy.exe        # Intérprete
└── lib/
    └── ...               # Librerías de soporte

LisPy/
├── lispy.c              # Código fuente
├── lib/                 # Librerías Lisp
├── docs/                # Documentación
└── README.md
\end{verbatim}

